\documentclass[11pt]{article}
\usepackage{fontspec}
\usepackage[french]{babel}
\usepackage{amsmath}
\usepackage{amssymb}
\usepackage{amsthm}
\usepackage{mathrsfs}
\usepackage{tikz-cd}
\usepackage{enumitem}

\newcommand{\printedpage}[1]{\marginpar{\footnotesize p.~#1}}
\newcommand{\msp}[1]{\marginpar{\footnotesize [#1]}}
\newcommand{\Parf}{\underline{\mathrm{Parf}}}
\newcommand{\Fib}{\underline{\mathrm{Fib}}}
\newcommand{\Liss}{\underline{\mathrm{Liss}}}
\newcommand{\Prop}{\underline{\mathrm{Prop}}}
\newcommand{\Bifib}{\underline{\mathrm{Bifib}}}
\newcommand{\Cofibsc}{\underline{\mathrm{Cofibsc}}}
\newcommand{\Fibsc}{\underline{\mathrm{Fibsc}}}
\newcommand{\HOT}{\mathrm{HOT}}
\newcommand{\margbar}{\marginpar{\(\|\)}}

\begin{document}

\printedpage{114}
% head: A. Grothendieck : Dérivateurs, Ch. VI, Édition des Universités de Montpellier II et Paris VII
% foot: Version 9 novembre 2001

$\Parf_\omega S$? Finalement, il me semble (contrairement à ce que j'avais cru voir) qu'il ne reste presque rien (\footnote{\textbf{NB} Si $W$ satisfait l'axiome des limites, alors W(4b$'$) implique W(4b), cf. théorème récapitulatif dans XVI 4.4.13 (p. 58-62), notamment $2^{o}$). La condition W(4b) est la condition B de cet énoncé, et W(4b$'$) la condition D. (Elle est impliquée par W(4a) $=$ C, ce que je soupçonnais nullement en écrivant XII \dots)}). On ne sait alors toujours pas, pour $X$ dans $\Fib_W S$, si
\[ i_X : X \longrightarrow C(X) \]
n'est dans $W^{\mathrm{f}}_S$, et non seulement dans $W_S$; on ne le sait que pour $X \in \Parf_\omega S$, mais dans ce cas on n'en a pas besoin, car l'utilité du foncteur $C(X)$ et de l'inclusion $i_X : X \longrightarrow C(X)$, c'est ‘d'améliorer’ $X$ et de le rendre parfait, voire même $W$-parfait. Mais s'il l'est déjà, c'est plus la peine. Il y aura, comme
\msp{page 184}
avant, 4 localisées essentiellement distinctes $\HOT_W$, $\HOT^{o}_W$, $\HOT^{\mathrm{lc}}_W$, $\HOT^{\mathrm{lc}\,!}_W$, plus une infinité d'autres $\mathcal{F}(W^{\mathrm{f}}_S)^{-1}$ pour $\mathcal{F}$ entre $\Parf_\omega S$ et $\Fib_W S$ (cf. page 180), avec le seul bénéfice qu'on peut maintenant exclure le cas $\mathcal{F} = \Parf_\omega S$, car dans ce cas $\mathcal{F}(W^{\mathrm{f}}_S)^{-1} = \mathcal{F}(W_S)^{-1} = \HOT^{\mathrm{lc}\,!}_W(S)$ fait partie des quatre canoniques. La chose décisive, c'est qu'on n'a toujours aucun moyen de prouver que
\[ \lambda : \HOT^{\mathrm{lc}}_W(S) \longrightarrow \HOT^{\mathrm{lc}\,!}_W(S) \]
est une équivalence, et je doute que ce soit le cas. Il faudrait regarder avec $W = W_{\mathbf{D}}$ la $\mathbf{D}$-équivalence, où $\mathbf{D}$ est un dérivateur Der($k$-Mod), $k$ un anneau donné (p.ex. $k = \mathbf{Z}$, ou $k$ un corps $\mathbf{Z}/p\mathbf{Z}$ ou $\mathbf{Q}$ \dots). Je ne vais pas le faire maintenant cependant. (Il est vrai que l'axiome W(4b$'$) n'est pas non plus satisfait dans ce cas \dots)

Si dans l'énoncé de l'axiome, on remplaçait $\Parf_\omega S$ par $\Parf_W S$, rien d'essentiel ne serait changé - tout au plus je remplacerais $\Parf_\omega S$ par $\Parf_W S$ dans le brillant diagramme p. 175. Finalement, je suis à présent convaincu qu'il faut tout ou rien : ou bien \textit{quatre} catégories $\HOT(S)$ pour $W$, en se contentant de W$(1,2,3)$, ou bien rien que trois, i.e. pas de $\HOT^{\mathrm{lc}\,!}_W(S)$ - et pour ça, il ne semble pas qu'on puisse se contenter de moins de W(4b).

\msp{page 185}
Il suffirait, il est vrai, de supposer que $W_S = W^{\mathrm{f}}_S$ sur l'une quelconque des sept catégories de l'étage autres que $\Parf_\omega S$. Cela impliquerait alors que dans les six autres, on aura $\overline{W_S} = \overline{W^{\mathrm{f}}_S}$ (mêmes saturés forts dans $\mathrm{Fl}(\mathcal{F})$). Un des deux les plus faibles des axiomes en question serait : % ?? uncertain: line under "étage" -- read as the overline of $\overline{W^{\mathrm{f}}_S}$ on the next line, not an underline of "étage"

\begin{quote}
Si $f : X \longrightarrow Y$ est un foncteur \textit{cartésien} de catégories Cat-fibrées et $W$-fibrées (\footnote{‘$W$-fibrées’ signifiant ici que pour $u \in \mathrm{Fl}\,S$, les foncteurs $u^*$ de changement de base $X_{s'} \longrightarrow X_s$, $Y_{s'} \longrightarrow Y_s$ sont dans $W$.}) sur une autre $S$, si $f \in W$, alors $f_s \in W$ pour tout $s \in S$,
\end{quote}

(et l'autre est la version duale, pour les foncteurs cocartésiens de catégories cofibrées). Sauf erreur, cet axiome doit déjà impliquer que $\overline{W}$ satisfait la version W(4a) la plus forte
\printedpage{115}
% head: A. Grothendieck : Dérivateurs, Ch. VI, Édition des Universités de Montpellier II et Paris VII
% foot: Version 9 novembre 2001
de toutes (a fortiori W(4b)) des axiomes W(4). Cet axiome me paraît donc à présent beaucoup plus raisonnable que W(4b$'$), et ceci d'autant plus qu'il se réfère directemment à la description de $\HOT^{\mathrm{lc}}_W(S)$ en termes de foncteurs (ou de ‘pseudofoncteurs’, si on ne veut supposer les catégories scindées et les foncteurs compatibles aux scindages) $S^{o} \longrightarrow \mathrm{Cat}$. % [sic] directemment -> directement, spelling
D'ailleurs, bien qu'il n'en ait pas l'air, cet axiome \textit{est} autodual, puisque la version duale, pour $S$ fixée, est la version directe pour $S^{o}$! Cela rend cet axiome encore plus appétissant. Donc pour formuler W(4a$'$), j'ai envie à présent de laisser tomber $\Parf_\omega$, et de faire pareil. Si on garde deux versions W(4a$'$) (l'ancienne) et W(4a$''$) (la nouvelle), l'une et l'autre ont la vertu d'impliquer W(4b). Et W(4b$'$) (nouvelle version) doit impliquer W(4a) (version ultraforte), sinon pour $W$, du moins pour $\overline{W}$.

Mes idées étaient tenacement [?] embrouillées
\msp{page 186}
jusqu'à aujourd'hui même, car j'avais l'idée que, grâce au mirifique foncteur $C$, la catégorie $\Parf_\omega S$, comme les sept autres catégories de l'étage, pourrait servir à définir $\HOT^{\mathrm{lc}}_W(S)$. Or il n'en est rien, semble-t-il, à moins de supposeer l'axiome W(4b), ou la variante W(4b$'$) (nouvelle version). L'ancienne version de W(4b$'$) s'avère parfaitement inutilisable. % [sic] supposeer -> supposer, spelling

À ce propos, je n'ai toujours pas répondu clairement pourtant, pas un ‘oui’ ou un ‘non’, à la question :

\begin{quote}
\textbf{Question :} \textit{les inclusions}
\[
\begin{tikzcd}
 & \Liss_{0,W}S \\
\Parf_W S \arrow[ur] \arrow[dr] & \\
 & \Prop_{0,W}S
\end{tikzcd}
\]
\textit{induisent-elles, même en l'absence de W(4b) (ou de ses variantes affaiblies) des équivalences pour les localisées fines (par $W^{\mathrm{f}}_S$)? Même question pour} $\Parf_\omega S$ \textit{au lieu de} $\Parf_W S$ (\footnote{OK en présence de l'axiome W(6bis) sur les $\varinjlim$ , cf. p. 192 ff.}). \textit{Même question pour}
\[ \Bifib\, S \subseteq \Parf_\omega S, \]
\textit{la catégorie des catégories sur $S$ qui sont à la fois} Cat-\textit{fibrées et} Cat-\textit{cofibrées.}
\end{quote}

Je ne me rends pas compte, même en admettant (cette fois) l'axiome W(4a), si $\Bifib\, S$ est assez grosse pour pouvoir décrire tous le types d'homotopie relatifs localement constants sur $S$, i.e. si % [sic] tous le types -> tous les types, missing letter
\[ (\Bifib\, S)(W^{\mathrm{f}}_S)^{-1} \longrightarrow (\Fib_W S)(W^{\mathrm{f}}_S)^{-1} \]
est au moins essentiellement surjectif, sinon une équivalence de catégories. Il manque un foncteur
\[ \mathrm{Cat}/S \longrightarrow \Bifib\, S \qquad \text{ou du moins} \qquad \underbrace{\Fib_W S}_{\text{ou $\Fib_\omega S$\,?}} \longrightarrow \Bifib\, S, \]
\printedpage{116}
% head: A. Grothendieck : Dérivateurs, Ch. VI, Édition des Universités de Montpellier II et Paris VII
% foot: Version 9 novembre 2001
qui puisse jouer le rôle de $\Phi$ et de $\Psi$ \textit{simultané\msp{page 187}ment.} Je n'ai pas même vérifié que $\Phi$ ne transforme \textit{pas} objets de $\Prop_W S$ ou de $\Prop_{0,W}(S)$ en objets $W$-propres (donc $W$-parfaits, étant de toutes façons $W$-lisse et même $W_\omega$-lisse) sur $S$ (\footnote{Voir réponse à cette intéressante question page 192.}). Si c'était pourtant vrai sur $\Prop_W S$, ça montrerait que l'image essentielle de $\Prop_W S$ dans $\HOT_W(S)$ ($= (\mathrm{Cat}/S)(W^{\mathrm{d}}_S)^{-1}$) est contenu dans $\HOT^{\mathrm{lc}}_W(S)$. % [sic] contenu -> contenue, agreement with "l'image"
Prenant $S = \mathbf{\Delta}^1$, cela signifierait que si dans $X$ on a une sous-catégorie fermée $X_1$ telle que $X_1 \overset{i}{\hookrightarrow} X$ soit $W_\omega$-asphérique, alors $1\backslash X \longrightarrow 0\backslash X$ est dans $W$ - or $0\backslash X \simeq X$, $1\backslash X \simeq X_1$, la flèche est $X_1 \longrightarrow X$, qui est bien dans $W$ étant $W$-asphérique! C'est étrange et quasiment immoral (\footnote{On observe plus bas que c'est ainsi chaque fois que $S$ a un objet final.}). Mais à force de gribouillis je viens de me convaincre que même si $X$ est cofibrée sur $S$ (donc aussi $W$-propre), $\Phi(X)$ n'est pas $W$-propre (ce qui impliquerait $W$-parfait), ni
\msp{page 188}
même un $W$-fibré, ce qui signifierait que
\[ s'\backslash X \longrightarrow s\backslash X \qquad \text{associé à } s \longrightarrow s' \text{ dans } S \]
serait toujours dans $W$. Posant $Z' = s'\backslash S$, $Z = s\backslash S$, cela signifierait que si on a un \textit{co-isomorphisme local}
\[ i : \underbrace{Z'}_{\simeq z_0\backslash Z} \longrightarrow Z, \]
où $Z'$, $Z$ ont des éléments initiaux, et si $X$ est cofibrée sur $Z$, $X'$ son image inverse sur $Z'$, alors $X' \longrightarrow X$ est dans $W$. En fait, si $Z$ a aussi un élément \textit{final} $e$ (comme c'était le cas tantôt pour $\mathbf{\Delta}^1$), alors $Z' \simeq z_0\backslash Z$ en a aussi en $e'$, qui s'envoie sur $e$, et on aura le diagramme commutatif
\[
\begin{tikzcd}
X' \arrow[r, "j"] & X \\
X'_{e'} \arrow[u, "W"] \arrow[r, "\sim"] & X_e, \arrow[u, "W"']
\end{tikzcd}
\]
d'où $j \in W$. Donc pour faire un contrexemple,
\msp{page 189}
il faut prendre $Z$ sans élément final.

Mais prenons $Z$ \texttt{[un]} ensemble ordonné, avec un objet initial $0$, mais sans objet final, et prenons pour $z_0$ un élément \textit{maximal} de $Z$, de sorte que $Z' = z_0\backslash Z$ soit isomorphe à $\{z_0\}$. Il faut donner dans Cat un exemple où l'inclusion de $X_{z_0}$ dans $X$ n'est pas dans $W$. On va prendre
\printedpage{117}
% head: A. Grothendieck : Dérivateurs, Ch. VI, Édition des Universités de Montpellier II et Paris VII
% foot: Version 9 novembre 2001
\[
\begin{tikzcd}
 & & a \\
Z = & 0 \arrow[ur] \arrow[dr] & \\
 & & b = z_0,
\end{tikzcd}
\]
donc la donnée de $X$ cofibrée coscindée sur $Z$ équivaut à un diagramme dans Cat
\[
\begin{tikzcd}
 & X_a \\
X_0 \arrow[ur, "\alpha"] \arrow[dr, "\beta"] & \\
 & X_b.
\end{tikzcd}
\]
Je pense que sous des conditions assez fortes (p.ex. si $\alpha$, $\beta$ sont des immersions ouvertes, ou des immersions fermées - à vérifier!), la catégorie $X$ est isomorphe dans Hot à la somme amalgamée dans Cat
\[ X_a{\textstyle\coprod_{X_0}}X_b, \]
et il n'y \texttt{[a]} aucune raison (sauf si on suppose $\alpha \in W$, et encore \dots) que $X_b \longrightarrow X_a \coprod_{X_0} X_b$ soit dans $W$. Prenons $X_0 = X_b = e$, $\beta = \mathrm{id}_e$, donc $\alpha$ est donné par la donnée d'un élément $t$ de $\mathrm{Ob}(X_a = T)$. La somme amalgamée est $T$, qui n'est pas isomorphe à $e$ dans $\mathrm{Hot}_W$,
\msp{page 190}
si $T$ n'a pas été prise $W$-asphérique. Mais je vais essayer de montrer directement que l'inclusion
\[ T = X_a \longrightarrow X \]
est $\in W$. Pour ceci je regarde la sous-catégorie ouverte $Z''$ de $Z$, complémentaire de la sous-catégorie pleine fermée $Z'$, et je note que $X'' = X|Z''$ étant propre sur $Z''$ et $a$ étant objet final de $Z''$, on a
\[ X_a \longrightarrow X'' \qquad \text{dans } W, \]
il reste à voir qu'il en est de même de $X'' \longrightarrow X$. Or $X''$ est déduite de $X$ en ajoutant un unique élément $\xi$ (celui de $X_b$), lequel ne s'envoie dans aucun élément de $X''$, et qui ne reçoit aucun élément de $X_a$ - par contre pour l'unique élément $\xi_0$ de $X_0$ il y a une unique flèche de celui-ci dans $\xi$. On a donc une rétraction de $X$ sur $X''$, envoyant $\xi$ sur $\xi_0$ et la flèche précédente sur $\mathrm{id}_{\xi_0}$. Est-ce un homotopisme? Il suffit de \texttt{[montrer]} qu'en tant qu'endomorphisme $\varphi$ de $X$, c'est $\sim \mathrm{id}_X$. Or on a bien $\mathrm{id}_X \longleftarrow \varphi$, dont la valeur sur $X''$ est l'identité, et la valeur sur $\xi$ est l'unique $\varphi(\xi) = \xi_0 \longrightarrow \xi$. C'est un peu bref, mais ça doit être vrai. (De toutes façons, $X$ est la réunion de deux ouverts $X''$, et de l'ouvert $X|\{0,b\} = X'''$, avec $X''' \simeq \mathbf{\Delta}^1$, et $X'' \cap X''' \simeq e$, donc au moins pour la cohomologie à coefficients localement constants de $X$ je suis à l'aise par Mayer-Vietoris, et sais bien que c'est celle du morceau $X''$, donc celle de $X_0$ \texttt{[ou $X_a$, indice illisible]} nullement celle de $e$!)

\printedpage{118}
% head: A. Grothendieck : Dérivateurs, Ch. VI, Édition des Universités de Montpellier II et Paris VII
% foot: Version 9 novembre 2001
\msp{page 191}
\margbar On a donc vu au moins ceci : Si $X$ est cofibrée donc $W$-propre sur $S$, il ne s'ensuit nullement que $\Phi(X)$ soit encore $W$-propre sur $S$ (\footnote{voir résultat plus précis page suivante!}), ni que ce soit un $W$-fibré. Mais que se passe-t-il si $X \in \Prop_{0,W}S$, par exemple s'il provient de $\Cofibsc_0 S$? Alors $X$ est dans $\Fib_W S$, donc on sait déjà que $\Phi(X)$ est dans $\Fibsc_{0,W}S$, donc ce n'est pas par là qu'on fera un contrexemple. Mais considérons encore $X$ cofibrée sur
\[
\begin{tikzcd}
 & & a \\
S = & 0 \arrow[ur] \arrow[dr] & \\
 & & b,
\end{tikzcd}
\]
donnée par le diagramme
\[
\begin{tikzcd}
 & X_a \\
X_0 \arrow[ur, "u"] \arrow[dr, "v"] & \\
 & X_b
\end{tikzcd}
\]
dans Cat, en supposant cette fois $u, v \in W$ ($X \in \Cofibsc_{0,W}S$). Alors $\Phi(X)$ est la catégorie fibrée scindée sur $S$ définie par le diagramme contravariant sur $S$
\[
\begin{tikzcd}
 & X_a \arrow[dl, hook, "i"'] \\
X_0 & \\
 & X_b, \arrow[ul, hook, "j"']
\end{tikzcd}
\]
où $i$, $j$ sont les inclusions (qui sont des immersions fermées). Dire que $\mathcal{X} = \Phi(X)$ est propre sur $S$ signifie aussi que ses restrictions au dessus des sous-catégories ouvertes $\{0,a\}$ et $\{0,b\}$ de $S$ le sont (car toute flèche $\mathbf{\Delta}^1 \longrightarrow S$ se factorise par l'une ou l'autre). Cela
\msp{page 192}
signifie donc aussi que l'inclusion des fibres fermées
\[ X_a \longrightarrow X, \qquad X_b \longrightarrow X \]
sont coasphériques. Mais $X_a \longrightarrow X$ ne peut être coasphérique que si $X_b = \emptyset$, car pour $x \in X_b \subseteq X$, $x\backslash X_a$ (ainsi que $X_a/x$) est vide, donc \textit{pas} $W$-asphérique (si on ne suppose $W = \mathrm{Fl}(\mathrm{Cat})$, cas que nous excluons). Donc $\Phi(X)$ n'est propre sur $S$ \textit{que} si $X_a$ et $X_b = \emptyset$ (résultat valable pour \textit{toute} $X$ sur $S$, qu'elle soit cofibrée sur $S$ ou $W$-propre, ou non). Si on veut $u, v \in W$, cela implique donc que $X_0 = \emptyset$, donc on a prouvé :

\end{document}
